
\section{Conclusion}

Continuity argues for a small but durable shift in how coding agents work with repositories. The repository should not store code alone. It should also store the minimal state needed to let work survive a session boundary: current truth, active tasks, validated lessons, append-only history, and reviewable evidence.

That claim is strongest when it remains narrow. External memory research supports persistent state. Repository-local manifest practice shows that current tools already lean this way. Architecture and provenance research explain why the state should be inspectable. Security research explains why the same files cannot be trusted blindly. The resulting design stance is clear: keep the state layer compact, scoped, validated, and easy to delete. Anything more should earn its place by evidence.
